\lecture{42}{Введение в федеративное обучение}

\sourcevideo
    {https://www.youtube.com/playlist?list=PLOzRYVm0a65fhKDS8cYIC7YCuVGJ4FOJx}
    {Week 11: Lecture 42: Introduction to Federated Learning}

Федеративное обучение предназначено для обучения общей модели по
данным, распределённым между клиентами. Исходные данные остаются на
устройствах или в организациях, а сервер получает только результаты
локальных вычислений. Клиентами могут быть мобильные телефоны,
устройства интернета вещей, больницы, банки и другие организации с
закрытыми наборами данных.

\section{Постановка задачи}

Пусть имеется \(K\) клиентов. Клиент \(k\) хранит набор
\(\mathcal D_k\) мощности \(n_k=\lvert\mathcal D_k\rvert\), а общее
число объектов равно
\[
    n=\sum_{k=1}^{K}n_k.
\]
Для функции потерь \(\ell(w;\zeta)\) локальная и глобальная функции
задаются формулами
\[
    F_k(w)
    =
    \frac1{n_k}
    \sum_{\zeta\in\mathcal D_k}
    \ell(w;\zeta),
    \qquad
    F(w)
    =
    \sum_{k=1}^{K}\frac{n_k}{n}F_k(w).
\]
Задача обучения имеет вид
\[
    \min_{w\in\mathbb R^d}F(w).
\]
Вес \(n_k/n\) учитывает объём локальных данных.

Типичный пример --- предсказание следующего слова на мобильной
клавиатуре. История ввода создаётся пользователем, может содержать
чувствительную информацию и существенно различается между
устройствами. Централизованный сбор таких данных увеличивает
требования к хранению, передаче и защите информации.

\section{Коммуникационный раунд}

Обычный сервер параметров на каждом шаге рассылает модель, получает
локальные градиенты и сразу обновляет параметры. При миллионах
клиентов эта схема требует высокой пропускной способности и
одновременной доступности большого числа устройств. Кроме того,
клиенты различаются по скорости вычислений, памяти, заряду батареи и
качеству соединения, а передаваемые обновления могут раскрывать
свойства локальных данных.

В федеративной схеме сервер на раунде \(t\) выбирает подмножество
\[
    S_t\subseteq\{1,\ldots,K\}.
\]
Если в раунде участвует доля \(C\in(0,1]\) клиентов, обычно берут
\[
    m=\max\{1,\lfloor CK\rfloor\}.
\]
Сервер отправляет выбранным клиентам текущую модель \(w^t\), клиенты
обучают её на своих данных и возвращают локальные модели или их
изменения. Наборы \(\mathcal D_k\) устройство не покидают.

Частичное участие необходимо, поскольку устройство может быть
отключено, занято пользователем, недостаточно заряжено или связано с
сервером медленным каналом. В практических системах локальное обучение
часто запускают только при выполнении заданных условий, например при
подключении телефона к питанию и беспроводной сети.

\section{Локальное обучение и агрегирование}

Получив \(w^t\), клиент \(k\) полагает
\[
    w_k^{t,0}=w^t
\]
и выполняет стохастические градиентные шаги
\[
    w_k^{t,s+1}
    =
    w_k^{t,s}
    -
    \eta\,
    g_k\bigl(w_k^{t,s};\mathcal B_k^{t,s}\bigr),
\]
где \(\eta>0\) --- локальный шаг, а
\(\mathcal B_k^{t,s}\subseteq\mathcal D_k\) --- мини-пакет. После
\(\tau_k\) шагов клиент отправляет серверу \(w_k^{t,\tau_k}\) или
приращение
\[
    \Delta_k^t=w_k^{t,\tau_k}-w^t.
\]

Если размер мини-пакета равен \(B\), а клиент выполняет \(E\)
локальных эпох, то
\[
    \tau_k
    =
    E\left\lceil\frac{n_k}{B}\right\rceil.
\]
При делимости \(n_k\) на \(B\) получаем \(\tau_k=En_k/B\). Поэтому
даже при одинаковых \(E\) и \(B\) клиенты с различными объёмами данных
выполняют разное число обновлений.

После получения локальных моделей сервер вычисляет
\[
    w^{t+1}
    =
    \sum_{k\in S_t}
    \frac{n_k}{\sum_{j\in S_t}n_j}
    w_k^{t,\tau_k}.
\]
Эквивалентная запись через локальные приращения имеет вид
\[
    w^{t+1}
    =
    w^t
    +
    \sum_{k\in S_t}
    \frac{n_k}{\sum_{j\in S_t}n_j}
    \Delta_k^t.
\]
Такое усреднение присваивает больший вес клиентам с большим числом
объектов. Подробная схема FedAvg рассматривается в следующей лекции.

Несколько локальных шагов уменьшают число коммуникационных раундов,
но усиливают расхождение локальных моделей. В сервере параметров
клиент обычно передаёт градиент после одного шага, тогда как в
федеративной схеме \(\tau_k\) может быть значительно больше единицы.
Чем дольше клиент оптимизирует собственную функцию \(F_k\), тем сильнее
его траектория отклоняется от направления глобальной функции \(F\).

\section{Статистическая и системная неоднородность}

В федеративных задачах локальные наборы обычно не имеют одинакового
распределения. Пользователи вводят разные слова, используют разные
языки и обладают различными привычками, поэтому функции \(F_i\) и
\(F_j\) могут существенно отличаться. При большом числе локальных
шагов это различие создаёт дрейф локальной модели к минимизатору
\(F_k\), а не глобальной функции.

Системная неоднородность связана с различием процессоров, объёмов
памяти, скоростей соединения, уровней заряда и времени локального
обучения. Часть выбранных клиентов может не завершить вычисления к
моменту агрегации; такие участники образуют проблему отстающих
клиентов. Таким образом, параметры \(C\), \(B\), \(E\), \(\eta\) и
\(\tau_k\) одновременно определяют вычислительную нагрузку,
коммуникационную стоимость и величину локального дрейфа.

\section{Конфиденциальность}

Локальное хранение \(\mathcal D_k\) уменьшает необходимость передачи
исходных данных, но не обеспечивает приватность автоматически. По
градиенту или обновлению модели иногда можно восстановить свойства
локальной выборки. Поэтому федеративную схему дополняют защищённой
агрегацией, шифрованием, дифференциальной приватностью, ограничением
нормы и зашумлением обновлений. При защищённой агрегации сервер видит
сумму клиентских обновлений, но не каждое обновление отдельно. Эти
механизмы являются дополнительными компонентами, а не следствием
самого факта локального обучения.

\section{Основные архитектуры}

В cross-device постановке участвует очень большое число мобильных и
IoT-устройств. В одном раунде доступна лишь малая доля клиентов, связь
нестабильна, вычислительные возможности различаются, а локальные
наборы обычно малы и неоднородны.

Cross-silo постановка объединяет меньшее число организаций, например
больниц, банков или исследовательских центров. Такие клиенты чаще
имеют стабильное соединение и более мощную инфраструктуру, однако их
данные также остаются локальными и могут существенно различаться.
Обе постановки используют серверную агрегацию общей модели и требуют
отдельных механизмов защиты информации.

Федеративное обучение не следует отождествлять с децентрализованным
SGD. В децентрализованном методе центрального сервера нет, и агенты
обмениваются моделями с соседями по графу. В стандартной федеративной
схеме сервер выбирает клиентов, рассылает модель, собирает локальные
обновления и выполняет агрегацию. Её определяющими признаками являются
локальное хранение данных, частичное участие и несколько локальных
шагов между коммуникациями.
